% GENERATED FILE. Edit canonical lessons, then run npm run generate:books.
% canonical-source-hash: fnv1a64:d1c8d3592ff52c05
% canonical-lessons: AR-C508-tafh-jildi, AR-C508-hikka, AR-C508-hakka, AR-C508-iltihab, AR-C508-tawarrama

\chapter{Skin trouble}
\label{ch:ar-skin-trouble}
\hlchaptermodality{508}

\begin{chapteropening}
\textbf{By the end of this chapter:} \emph{I can talk about a rash, an itch and an inflammation, and say that someone scratched or that a foot swelled up.}

Complete the last lesson of chapter 508: I can talk about a rash, an itch and an inflammation, and say that someone scratched or that a foot swelled up.
\end{chapteropening}

\section[\texorpdfstring{\ar{طفح} \ar{جلدي} (ṭafḥ jildī)}{طفح جلدي (ṭafḥ jildī)}]{\ar{طفح} \ar{جلدي} (ṭafḥ jildī) — a rash}

\label{lesson:AR-C508-tafh-jildi}



\begin{quote}
Before the new one: say the Arabic for to vomit, then the Arabic for diarrhoea.
\end{quote}

\subsection*{You'll want to know: \ar{طفح} \ar{جلدي}}
\textbf{\ar{طفح} \ar{جلدي}} — \emph{ṭafḥ jildī} — \textquotedblleft{}a rash\textquotedblright{}.

\begin{quote}
\ar{ظهر} \ar{طفح} \ar{جلدي} \ar{على} \ar{يد} \ar{الطفل}.
\end{quote}

\begin{quote}
\emph{ẓahara ṭafḥ jildī ʿalā yad aṭ-ṭifl.}
\end{quote}

\begin{quote}
A rash appeared on the child's hand.
\end{quote}

Often shortened to \ar{طفح}.

\begin{cousinweb}
\textbf{\ar{طفح} \ar{جلدي}} (\emph{ṭafḥ jildī}) is built on two native Arabic roots: \emph{ṭafḥ} from \emph{ṭ-f-ḥ}, \textquotedblleft{}to fill to the brim, to overflow\textquotedblright{} (Lane), and \emph{jildī}, \textquotedblleft{}of the skin\textquotedblright{}, from \emph{jild}, skin, which you know. A rash is the skin brimming over.
\end{cousinweb}

\begin{grammarlens}[title={What you've built}]
One more word for describing what you lived and felt.
\end{grammarlens}

\subsection*{Your turn}
\begin{itemize}
\raggedright
  \item \emph{Say it:} \emph{ṭafḥ jildī}
  \item \emph{Say it:} \emph{ṭafḥ jildī}, once more
  \item \emph{Say it:} \emph{ẓahara ṭafḥ jildī ʿalā yad aṭ-ṭifl.} — A rash appeared on the child's hand.
  \item \emph{Recall:} say the Arabic for to vomit, then the Arabic for diarrhoea, then say \emph{ṭafḥ jildī} again
\end{itemize}

\begin{tcolorbox}[breakable,colback=teal!4,colframe=teal!35!black,title={Before you move on}]
What does \ar{طفح} \ar{جلدي} mean? (\textquotedblleft{}A rash\textquotedblright{}.) Say it once more.
\end{tcolorbox}
\section[\texorpdfstring{\ar{حكة} (ḥikka)}{حكة (ḥikka)}]{\ar{حكة} (ḥikka) — an itch, itching}

\label{lesson:AR-C508-hikka}



\begin{quote}
Before the new one: say the Arabic for diarrhoea, then the Arabic for a rash.
\end{quote}

\subsection*{You'll want to know: \ar{حكة}}
\textbf{\ar{حكة}} — \emph{ḥikka} — \textquotedblleft{}an itch, itching\textquotedblright{}.

\begin{quote}
\ar{الحكة} \ar{تزعجني} \ar{في} \ar{الليل}.
\end{quote}

\begin{quote}
\emph{al-ḥikka tuzʿijunī fī al-layl.}
\end{quote}

\begin{quote}
The itching bothers me at night.
\end{quote}

\begin{cousinweb}
\textbf{\ar{حكة}} (\emph{ḥikka}) is built on the native Arabic root \emph{ḥ-k-k}, \textquotedblleft{}to rub, to scratch\textquotedblright{} (Lane). Ḥikka is the itch that makes you scratch.
\end{cousinweb}

\begin{grammarlens}[title={What you've built}]
One more word for describing what you lived and felt.
\end{grammarlens}

\subsection*{Your turn}
\begin{itemize}
\raggedright
  \item \emph{Say it:} \emph{ḥikka}
  \item \emph{Say it:} \emph{ḥikka}, once more
  \item \emph{Say it:} \emph{al-ḥikka tuzʿijunī fī al-layl.} — The itching bothers me at night.
  \item \emph{Recall:} say the Arabic for diarrhoea, then the Arabic for a rash, then say \emph{ḥikka} again
\end{itemize}

\begin{tcolorbox}[breakable,colback=teal!4,colframe=teal!35!black,title={Before you move on}]
What does \ar{حكة} mean? (\textquotedblleft{}An itch, itching\textquotedblright{}.) Say it once more.
\end{tcolorbox}
\section[\texorpdfstring{\ar{حك} (ḥakka)}{حك (ḥakka)}]{\ar{حك} (ḥakka) — to scratch}

\label{lesson:AR-C508-hakka}



\begin{quote}
Before the new one: say the Arabic for a rash, then the Arabic for an itch.
\end{quote}

\subsection*{You'll want to know: \ar{حك}}
\textbf{\ar{حك}} — \emph{ḥakka} — \textquotedblleft{}to scratch\textquotedblright{}.

\begin{quote}
\ar{حك} \ar{الولد} \ar{رأسه} \ar{وفكر}.
\end{quote}

\begin{quote}
\emph{ḥakka al-walad raʾsahu wa-fakkara.}
\end{quote}

\begin{quote}
The boy scratched his head and thought.
\end{quote}

Present: \ar{يحك}, \emph{yaḥukk}.

\begin{cousinweb}
\textbf{\ar{حك}} (\emph{ḥakka}) is built on the native Arabic root \emph{ḥ-k-k}, \textquotedblleft{}to rub, to scratch\textquotedblright{} (Lane). Its last two root letters are the same, so they write as one doubled letter.
\end{cousinweb}

\begin{grammarlens}[title={What you've built}]
One more word for describing what you lived and felt.
\end{grammarlens}

\subsection*{Your turn}
\begin{itemize}
\raggedright
  \item \emph{Say it:} \emph{ḥakka}
  \item \emph{Say it:} \emph{ḥakka}, once more
  \item \emph{Say it:} \emph{ḥakka al-walad raʾsahu wa-fakkara.} — The boy scratched his head and thought.
  \item \emph{Recall:} say the Arabic for a rash, then the Arabic for an itch, then say \emph{ḥakka} again
\end{itemize}

\begin{tcolorbox}[breakable,colback=teal!4,colframe=teal!35!black,title={Before you move on}]
What does \ar{حك} mean? (\textquotedblleft{}To scratch\textquotedblright{}.) Say it once more.
\end{tcolorbox}
\section[\texorpdfstring{\ar{التهاب} (iltihāb)}{التهاب (iltihāb)}]{\ar{التهاب} (iltihāb) — an inflammation}

\label{lesson:AR-C508-iltihab}



\begin{quote}
Before the new one: say the Arabic for an itch, then the Arabic for to scratch.
\end{quote}

\subsection*{You'll want to know: \ar{التهاب}}
\textbf{\ar{التهاب}} — \emph{iltihāb} — \textquotedblleft{}an inflammation\textquotedblright{}.

\begin{quote}
\ar{عند} \ar{أختي} \ar{التهاب} \ar{في} \ar{الحلق}.
\end{quote}

\begin{quote}
\emph{ʿinda ukhtī iltihāb fī al-ḥalq.}
\end{quote}

\begin{quote}
My sister has a sore throat.
\end{quote}

\ar{التهاب} \ar{الحلق} is a sore throat.

\begin{cousinweb}
\textbf{\ar{التهاب}} (\emph{iltihāb}) is built on the native Arabic root \emph{l-h-b}, \textquotedblleft{}a flame, to blaze\textquotedblright{} (Lane). Iltihāb is the verbal noun of Form VIII, a flaring up: the body burning in one spot.
\end{cousinweb}

\begin{grammarlens}[title={What you've built}]
One more word for describing what you lived and felt.
\end{grammarlens}

\subsection*{Your turn}
\begin{itemize}
\raggedright
  \item \emph{Say it:} \emph{iltihāb}
  \item \emph{Say it:} \emph{iltihāb}, once more
  \item \emph{Say it:} \emph{ʿinda ukhtī iltihāb fī al-ḥalq.} — My sister has a sore throat.
  \item \emph{Recall:} say the Arabic for an itch, then the Arabic for to scratch, then say \emph{iltihāb} again
\end{itemize}

\begin{tcolorbox}[breakable,colback=teal!4,colframe=teal!35!black,title={Before you move on}]
What does \ar{التهاب} mean? (\textquotedblleft{}An inflammation\textquotedblright{}.) Say it once more.
\end{tcolorbox}
\section[\texorpdfstring{\ar{تورم} (tawarrama)}{تورم (tawarrama)}]{\ar{تورم} (tawarrama) — to swell (up)}

\label{lesson:AR-C508-tawarrama}



\begin{quote}
Before the new one: say the Arabic for to scratch, then the Arabic for an inflammation.
\end{quote}

\subsection*{You'll want to know: \ar{تورم}}
\textbf{\ar{تورم}} — \emph{tawarrama} — \textquotedblleft{}to swell (up)\textquotedblright{}.

\begin{quote}
\ar{تورم} \ar{كاحلي} \ar{بعد} \ar{السقوط}.
\end{quote}

\begin{quote}
\emph{tawarrama kāḥilī baʿda as-suqūṭ.}
\end{quote}

\begin{quote}
My ankle swelled up after the fall.
\end{quote}

Present: \ar{يتورم}, \emph{yatawarram}.

\begin{cousinweb}
\textbf{\ar{تورم}} (\emph{tawarrama}) is built on the native Arabic root \emph{w-r-m}, \textquotedblleft{}to swell\textquotedblright{} (Lane). Form V tawarrama is to become swollen.
\end{cousinweb}

\begin{grammarlens}[title={What you've built}]
One more word for describing what you lived and felt.
\end{grammarlens}

\subsection*{Your turn}
\begin{itemize}
\raggedright
  \item \emph{Say it:} \emph{tawarrama}
  \item \emph{Say it:} \emph{tawarrama}, once more
  \item \emph{Say it:} \emph{tawarrama kāḥilī baʿda as-suqūṭ.} — My ankle swelled up after the fall.
  \item \emph{Recall:} say the Arabic for to scratch, then the Arabic for an inflammation, then say \emph{tawarrama} again
\end{itemize}

\begin{tcolorbox}[breakable,colback=teal!4,colframe=teal!35!black,title={Before you move on}]
What does \ar{تورم} mean? (\textquotedblleft{}To swell (up)\textquotedblright{}.) Say it once more.
\end{tcolorbox}
